% ======================================================================
% 3.3 Динамика регулятора. Ответственный: B — Попадюк Евгений.
%
% Что сдать (main.tex: теория 1281–1315, порядок 1316–1335,
% разобранный пример 1337–1375; требования к отчёту 1871–1879):
%   1. Найти положение равновесия (q*, c*) системы модели.
%   2. Найти матрицу Якоби J правых частей в равновесии, обозначив
%      A=lambda*(1-u0)^2/u0; объяснить, почему d(dot c)/dc = 0.
%   3. Найти собственные числа nu_{1,2}; показать устойчивость при любом k>0;
%      найти порог k*, при превышении которого возникают колебания (флаппинг).
%   4. Для своего lambda вычислить A и k* при u0=0.6; 0.75; 0.9;
%      сделать вывод о поведении системы при u0 -> 1.
%
% Параметры варианта 3 (приложение Б): mu=80, lambda=600.
%
% Программные значения — task_3_3 (python experiments.py 3.3).
% Численный порог флаппинга — п. 3.5.4 (файл 353-simulator.tex, зона C):
% он сравнивается с k*, но расхождение НЕ называется ошибкой аппроксимации —
% результат зависит от критерия обнаружения (main.tex, строки 1652–1661).
% ======================================================================

\subsection{Динамика регулятора}

\taskhead{Исходные данные и положение равновесия}
\todo{B: найти $(q^*,c^*)$ системы модели; вычислить для заданных целевых загрузок (пример: main.tex, строки 1338–1346).}

\taskhead{Матрица Якоби в равновесии}
\todo{B: выписать частные производные правых частей; составить $J$ в равновесии с обозначением $A=\lambda(1-u_0)^2/u_0$; объяснить, почему $\partial\dot c/\partial c=0$ (пример: main.tex, строки 1347–1362).}

\taskhead{Собственные числа и устойчивость}
\todo{B: найти $\nu_{1,2}$; обосновать локальную асимптотическую устойчивость при любом $k>0$; вывести порог $k_*$ появления комплексных собственных чисел (колебаний); пояснить, что этот порог — не порог потери устойчивости.}

\taskhead{Расчёт для трёх целевых загрузок}
\todo{B: вычислить $A$ и $k_*$ при $u_0=0.6;\,0.75;\,0.9$ для своего $\lambda$; сделать вывод о поведении порога при $u_0\to1$.}
\begin{table}[htbp]
\centering\small
\caption{Равновесие и порог колебаний для трёх целевых загрузок}
\label{tab:3-3}
\begin{tabular}{|r|r|r|r|r|}\hline
$u_0$ & $q^*$ & $c^*$ & $A$ & $k_*$\\\hline
0.6 & — & — & — & —\\
0.75 & — & — & — & —\\
0.9 & — & — & — & —\\
\hline\end{tabular}
\end{table}

\taskhead{Сопоставление с программным расчётом}
% Программные значения — вызовы своих функций из solution.py (task_3_3).
\todo{B: листинг вызовов функций и запуск python experiments.py 3.3.}
\begin{table}[htbp]
\centering\footnotesize
\setlength{\tabcolsep}{4pt}
\caption{Сравнение аналитических и программных значений (п. 3.3)}
\label{tab:3-3-check}
\begin{tabular}{|l|r|r|r|r|}\hline
Величина & \shortstack{Аналитическое\\значение} & \shortstack{Программное\\значение} & \shortstack{Абсолютная\\ошибка} & \shortstack{Относительная\\ошибка}\\\hline
$k_*$ при $u_0=0.6$ & — & — & — & —\\
$k_*$ при $u_0=0.75$ & — & — & — & —\\
$k_*$ при $u_0=0.9$ & — & — & — & —\\
\hline\end{tabular}
\end{table}

\taskhead{Вывод}
\todo{B: устойчивость при любом $k>0$; смысл порога $k_*$; поведение при $u_0\to1$.}
